\chapter{Introduction}

How can one choose a meaningful scale from a multiscale family of statistical objects?

\section{Basics about Topology}

\begin{itemize}
    \item definition of continuous function in topological space 
    \item contractible sets, LS-category, examples, homotopic functions
    \item Closure, Star, Link
    \item topolgical invariants (homotopy invariance)
\end{itemize}

\begin{Def}[Topological Space]
    Let $X$ be any set. Let $\mathcal{T}\subseteq\mathcal{P}(X)$ be a collection of subsets of $X$ such that
    \begin{itemize}
        \item $\emptyset\in\mathcal{T}$ and $X\in\mathcal{T}$
        \item $\bigcup_{i\in I}U_i\in\mathcal{T}$ for all $(U_i)_{i\in I}\subseteq\mathcal{T}$
        \item $U_1\cap U_2\in\mathcal{T}$ for all $U_1,U_2\in\mathcal{T}$
    \end{itemize}
    Then the pair $(X,\mathcal{T})$ is called a topological space.
\end{Def}

\begin{Def}[contractible set]
    Let $(X,\mathcal{T})$ be a topological space. A set $U\in\mathcal{T}$ is contractible in $X$ if there exists a (continuous) map $H:U\times[0,1]\to X$ and a point $x_0\in X$ such that $H(u,0)=u$ and $H(u,1)=x_0$ for all $u\in U$.
\end{Def}

\begin{Def}[Lusternik-Schnirelmann category]
    Let $(X,\mathcal{T})$ be a topological space. The Lusternik-Schnirelman category (LS-category) of $X$ is the minimum number of open contractible sets $U_i$ that cover $X$. Formally speaking
    \begin{align*}
        \text{cat}(X)=\min\{n\in\mathbb{N}:\exists U_1,\dots,U_n\in\mathcal{T},U_j\ \text{is contractible in}\ X\}
    \end{align*}
\end{Def}

The LS-category of a topological space can be interpreted as a measure for its complexity, i.e. a topological space has a more complex structure if we need more open and contractible sets to cover it.

\begin{Prop}
    Let $X_1,\mathcal{T}_1$ and $(X_2,\mathcal{T}_2)$ be topological spaces. Then the following properties hold:
    \begin{enumerate}
        \item[1.] Subadditivity: $\text{cat}(X_1\cup X_2)\le\text{cat}(X_1)+\text{cat}(X_2)$
        \item[2.] Simple upper bound
        \item[3.] Simple lower bound?
        \item[4.] Homotopy invariance
        \item[5.] $\text{cat}(X)=1\iff X$ is contractible
    \end{enumerate}
\end{Prop}



\section{Simplicial Complexes}

We begin by introducing the simplicial complex, the main mathematical object in Topological Data Analysis. For literature consult \cite{Dey_Wang_2022}.

\begin{Def}[Abstract simplicial complex]
    Let $V$ be a finite set of vertices. An abstract simplicial complex $K$ is a collection of subsets of $V$ such that:
    \begin{itemize}
        \item $\emptyset\in K$
        \item $\{v\}\in K$ for every $v\in V$
        \item if $\sigma\in K$ and $\tau\subseteq\sigma$ then $\tau\in K$
    \end{itemize}
    We denote by $\mathcal{K}(V)$ the set of all simplicial complexes with vertex set $V$.
\end{Def}

An element $\sigma\in K$ is called face. The dimension of a face $\sigma$ is $\dim{\sigma}=|\sigma|-1$. The dimension of $K$ is $\dim{K}=\max\{\dim{\sigma}:\sigma\in K\}$. A face $\sigma\in K$ with $\dim{\sigma}=\dim{K}$ is called facet.

Note that a simplicial complex $K$ with $\dim{K}=1$ is a simple graph. Thus, an abstract simplicial complex can be considered as a generalization of a simple graph.

The abstract simplicial complex is a purely combinatorial object. It can be given a geometric meaning by embedding it in the $d$-dimensional euclidean space.

\begin{Def}[Geometric realization]
    The geometric realization of an abstract simplicial complex $K$ with vertex set $V=\{v_1,\dots,v_n\}$ is given by
    \begin{align*}
        |K|_{\Phi}=\bigcup\limits_{\sigma\in K}\operatorname{conv}\{p_t:t\in\sigma\},
    \end{align*}
    where $\Phi:V\to\mathbb{R}^m$ is an embedding map that may have further properties.
\end{Def}

\todo[inline]{Geometric volume depends on the embedding.}

\todo[inline]{$2d+1$ theorem}

Using the geometric realization $|K|_{\Phi}$ the volume of a simplex $\sigma=[v_{i_0}\dots v_{i_j}]$ can be defined by
\begin{align*}
    \operatorname{Vol}(\sigma):= \frac{1}{j!}\sqrt{\det{\left(\langle x_{i_k}-x_{i_0},x_{i_l}-x_{i_0}\rangle_{k,l=1}^{j}\right)}},
\end{align*}
where $\langle x_k,x_l\rangle_{k,l=1}^{j}$ is the Gram matrix defined by $X^TX$ with $X=(x_1\cdots x_j)$. We see that $\operatorname{Vol}(\sigma)=0$ for all $\sigma\in K$ with $\dim{\sigma}<\dim{K}$. The volume of a simplicial complex is
\begin{align*}
    \operatorname{Vol}(K)=\sum_{\sigma\in K,\dim{\sigma}=\dim{K}}\operatorname{Vol}(\sigma)
\end{align*}

\subsection*{Binary operations on simplicial complexes}

\todo[inline]{Symmetric difference}

There exist different operations $\oplus:\mathcal{K}(V)\times\mathcal{K}(V)\to\mathcal{K}(V)$ which we will discuss below. The union of two complexes $K_1$ and $K_2$ is defined as
\begin{align*}
    K_1\cup K_2:=\{\sigma\subseteq V:\sigma\in K_1\text{ or }\sigma\in K_2\}
\end{align*}

\todo[inline]{Give general discussion on operations $\oplus\in\{\cup,\dots\}$ with examples}

% The geometric realization $\tilde{K}$ of an abstract simplicial complex $K$ can be viewed as a topological space, which allows us to define the LS-category of a simplicial complex $K$ by $\text{cat}(K):=\text{cat}(\tilde{K})$.



\subsection{Simplex Trees}

Simplex trees are used to store simplicial complexes efficiently.

\subsection{Types of simplicial complexes}
In this chapter, we describe different methods to construct simplicial complexes from data points $X=\{x_1,\dots,x_n\}\subseteq\mathbb{R}^d$.

\subsubsection{\v{C}ech complex}

$\check{C}_{\epsilon}(X)=\{\sigma\subseteq X:\bigcap\limits_{x\in\sigma} B(x,\epsilon)\ne\emptyset\}$

\subsubsection{Vietoris--Rips Complex}

$\operatorname{VR}_{\epsilon}(X)=\{\sigma\subseteq X:d(x,y)\le\epsilon\ \text{for all} \ x,y\in\sigma\}$

In what way is the Vietoris--Rips complex an approximation of the \v{C}ech complex?

\subsubsection{alpha complex}

\subsubsection{(Lazy) Witness complex}

\section{Topological Data Analysis}

\begin{itemize}
    \item persistance diagrams
    \item betti numbers
    \item homology
\end{itemize}

\todo[inline]{Cite papers about TDA intro}

Let $(V,\Sigma)$ be an abstract simplicial complex and denote by $\Sigma_k$ the set of all simplices in $\Sigma$ with dimension $k$. Let $C_k(\mathbb{F})$ be the free Abelian group generated using $\Sigma_k$ as the basis with linear coefficients from the field $\mathbb{F}$, that is,
\begin{align*}
    C_k(\mathbb{F})=\left\{\sum_{i=1}^{n}k_i\sigma_i:\sigma_i\in\Sigma_K,k_i\in\mathbb{F},n\in\mathbb{N},i=1,\dots,n\right\}.
\end{align*}
Throughout this work, we will only consider $\mathbb{F}=\mathbb{Z}_2=\{0,1\}$ and set $C_k(\mathbb{Z}_2)=C_k$. The sum in [] is seen as a formal sum: the coefficient $k_i=0$ means that the simplex $\sigma_i$ is excluded and the coefficient $k_i=1$ means that the simplex $\sigma_i$ is included in the sum. Next we define the boundary map $\delta_k:\Sigma_k\to C_{k-1}$ by
\begin{align*}
    \delta_k(\sigma):=\sum_{i=0}^{k}(-1)^{i}\sigma_{-i},\quad\sigma=[x_0\dots x_k]\in\Sigma_k,
\end{align*}
where $\sigma_{-i}\in\Sigma_{k-1}$ is the simplicial complex with vertices $x_0,\dots,x_{i-1},x_{i+1},\dots,x_k$. We extend $\delta_k$ from $\Sigma_k$ to $C_k$ by the rule
\begin{align*}
    \delta_k(\Gamma):=\sum_{i=1}^{n}k_i\delta_k(\sigma_i),
\end{align*}
where $\Gamma=\sum_{i=1}^{n}k_i\sigma_i\in C_k$.

The $k$-th Betti number $\beta_k$ is set to be the dimension of the image of $\delta_k$:
\begin{align*}
    \beta_k=\dim(\operatorname{Im}\delta_k)
\end{align*}